% !PRINTED 211
\weberpage{211}
By transformation by means of the substitution

\[
\begin{pmatrix} e^{\frac{\pi i}{4}}, & 0 \\ 0, & e^{-\frac{\pi i}{4}} \end{pmatrix}
\]

the dihedral group $D_m$ passes over into

\[
\begin{pmatrix} \varepsilon^{1/2\,\lambda}, & 0 \\ 0, & \varepsilon^{-1/2\,\lambda} \end{pmatrix}, \quad
\begin{pmatrix} 0, & \varepsilon^{1/2\,\lambda} \\ -\varepsilon^{-1/2\,\lambda}, & 0 \end{pmatrix}
\]

or into

\[
\begin{gathered}
x, \; \varepsilon x, \; \varepsilon^2 x, \; \ldots, \; \varepsilon^{m-1} x\\[4pt]
\dfrac{-1}{x}, \; \dfrac{-\varepsilon}{x}, \; \dfrac{-\varepsilon^2}{x}, \; \cdots, \; \dfrac{-\varepsilon^{m-1}}{x}
\end{gathered}
\]

and the fundamental forms for this representation are:

\[
f_1 = x_1 x_2, \quad f_2 = x_1^m + i^m x_2^m, \quad f_3 = x_1^m - i^m x_2^m.
\]

The dihedral group $D_m$ contains as divisor the cyclic
group $C_m$, and indeed $C_m$ is a \gesperrt{normal divisor} of $D_m$.

By transformation with any multiplicative
substitution

\[
\begin{pmatrix} \alpha, & 0 \\ 0, & \delta \end{pmatrix}
\]

$C_m$ does not change, while the substitutions of $D_m$ not contained
in $C_m$ change their form; for we have:

\[
\begin{gathered}
\begin{pmatrix} \delta, & 0 \\ 0, & \alpha \end{pmatrix}
\begin{pmatrix} \varepsilon, & 0 \\ 0, & 1 \end{pmatrix}
\begin{pmatrix} \alpha, & 0 \\ 0, & \delta \end{pmatrix}
=
\begin{pmatrix} \varepsilon, & 0 \\ 0, & 1 \end{pmatrix},\\[6pt]
\begin{pmatrix} \delta, & 0 \\ 0, & \alpha \end{pmatrix}
\begin{pmatrix} 0, & \varepsilon \\ 1, & 0 \end{pmatrix}
\begin{pmatrix} \alpha, & 0 \\ 0, & \delta \end{pmatrix}
=
\begin{pmatrix} 0, & \varepsilon \delta^2 \\ \alpha^2, & 0 \end{pmatrix}.
\end{gathered}
\]

Apart from such a multiplicative substitution, however,
the group $D_m$ is completely determined by the group $C_m$
contained in it. For if we set

\[
D_m = C_m + \varphi C_m,
\]
in which

\[
\varphi = \begin{pmatrix} \alpha, & \beta \\ \gamma, & \delta \end{pmatrix}
\]

is put, then, since $C_m$ must be a normal divisor of $D_m$, we must
have $\varphi C_m = C_m \varphi$, i.\,e.\ for every exponent $r$ there must
be an exponent $s$, and conversely, such that

\[
\begin{pmatrix} \alpha, & \beta \\ \gamma, & \delta \end{pmatrix}
\begin{pmatrix} \varepsilon^r, & 0 \\ 0, & 1 \end{pmatrix}
=
\begin{pmatrix} \varepsilon^s, & 0 \\ 0, & 1 \end{pmatrix}
\begin{pmatrix} \alpha, & \beta \\ \gamma, & \delta \end{pmatrix}
\]
%CONT>>
